\section{Two compact realizations of the reciprocal recorder}
\label{subsec:p2:rapid-periodic}
The preceding area-changing constructions use the third coordinate to
compensate a planar determinant. For the ordinary one-head recorder,
the planar maps already have determinant one. The third coordinate can
then serve only to separate instructions during transport. We give two
compact Euclidean realizations: one uses four stationary fields in
succession, and the other follows entire moving boxes in three phases.
Their observer estimates will follow directly from the horizontal
motion, without a vertical restriction on the detector.

Use the frontier recorder of Lemma~\ref{lem:sb-frontier} throughout this
section. Its frontier starts at absolute cell two; every other history
cell is blank and every return mark is zero. Its fixed incoming moves
and distinguishing written symbols will separate the target rectangles.
At checkpoint $k$ the records occupy $2,\ldots,k+1$, the frontier is at
$k+2$, and transition $k+1$ takes $2(k+2-h_{k+1})+4$ instructions.
These properties hold for the same entire partial table as in that lemma.

Choose the combined alphabet $\mathcal A$ of the recorder, of size $m$.
Give its letters the even digits $g_a=0,2,\ldots,2m-2$, with the fully
blank combined letter first. We will use either $B=2m$ or $B=2m+1$.
For a tape $v$ whose head is at its relative index zero, set
\[
 x=\sum_{j\ge1}g_{v_{-j}}B^{-j},\qquad
 y=\sum_{j\ge0}g_{v_j}B^{-j-1},\qquad
 I_a=[g_a/B,(g_a+1)/B].
\]
These are the gapped codes of Section~\ref{subsec:sb-codes}.
A state $s$ is placed at $(o_s+x,y,0)$ for a specified offset $o_s$.
If a rule reads $a$, writes $b$, moves by $d$, and enters $s'$, its
coordinate map relative to the source and target offsets is
\begin{equation}\label{eq:p2:rapid-rectangle-rules}
\begin{array}{c|c|c}
 d&\text{source rectangle}&(x',y')\\\hline
 1 &[0,1]\times I_a&((g_b+x)/B,\ By-g_a)\\
 0 &[0,1]\times I_a&(x,\ y+(g_b-g_a)/B)\\
 -1&I_\ell\times I_a&
 (Bx-g_\ell,\ (g_\ell+(g_b+By-g_a)/B)/B).
\end{array}
\end{equation}
For the last row there is one branch for each $\ell\in\mathcal A$.
The target rectangles are respectively $I_b\times[0,1]$,
$[0,1]\times I_b$, and $[0,1]\times(g_\ell/B+I_b/B)$.
These formulas follow by removing the scanned digit and pushing the
written digit. They apply to every point of the displayed closed
rectangles. Determinism separates the sources; the fixed incoming
displacement and distinguishing written symbol separate the targets.
For a left move the target pair $(\ell,b)$ supplies that distinction.
Every linear part is $\operatorname{diag}(\lambda,\lambda^{-1})$.

The code also retains absolute, rather than only head-relative,
configuration information at the machine checkpoints. Count the records
immediately to the left of the frontier to recover $k$. The frontier
has absolute position $k+2$; its location in the decoded tape therefore
recovers the absolute head position. The work track then recovers the
original tape. This use of history records is the retained-history
version of the reversible-computation idea of Bennett~\cite{Bennett1973};
the positional affine representation follows the generalized-shift
viewpoint of Moore~\cite{Moore1990,Moore1991}.

\begin{proposition}[Two compact periodic processors]
\label{prop:p2:rapid-periodic}
For every deterministic machine and finite input, and every fixed
computable viscosity $\nu>0$, there are effective smooth velocities
$U$ and forces $f$ on $\mathbb R^3$ with the following properties.
Both have all mixed derivatives bounded, support in a compact spatial
set independent of time, and period one for $t\ge1$; $U(0)=0$.
The solution with pressure zero is unique in the class
\[
 u\in C([0,T];H^2)\cap C^1([0,T];L^2),\qquad
 p\in C([0,T];H^1),\qquad
 \sup_{[0,T]\times\mathbb R^3}(|u|+|\nabla u|)<\infty
\]
for every finite $T$, and its kinetic energy is uniformly bounded.
The following are two separate constructions with their stated fixed
particle events:
\begin{enumerate}
\item Four successive stationary localized fields, with initial label
$(-1,0,0)$ and detector $\{x_1>0\}$.
\item Three successive phases of moving full-box fields, with initial
label $(0,0,0)$ and detector \mbox{$\{x_1<-1\}$}.
\end{enumerate}
In each case the particle enters the detector at some real time if and
only if the machine halts. The repeated part of the prescription depends
only on the machine; the input appears only in a one-time loading pulse.
\end{proposition}
\begin{proof}
Let $R_i,R_i'$ be the planar source and target rectangles in
\eqref{eq:p2:rapid-rectangle-rules}, with centers $c_i,c_i'$ and first
scale $\lambda_i$. All their side lengths are at most one. If there are
no instructions, use the zero repeated field and the same loader described
below. Empty pairwise margin tests are omitted for a singleton family.

For the first realization choose $B=2m$. Put terminal offsets at
$4,7,10,\ldots$ and all other offsets at $-4,-7,-10,\ldots$.
Choose private heights $h_i=4i$. On four successive quarters of a unit
period, perform the following autonomous motions, each with flat progress
from zero to one:
\begin{enumerate}
\item Translate upward by $h_i e_3$, with a cutoff equal to one near
$R_i\times[0,h_i]$.
\item At height $h_i$, use the Hamiltonian vector potential
\[
 (0,0,\chi_i\log(\lambda_i)
                (X_1-c_{i1})(X_2-c_{i2})).
\]
Its curl on the plateau is
$\log(\lambda_i)(X_1-c_{i1},-(X_2-c_{i2}),0)$.
\item Translate horizontally by $c_i'-c_i$ at height $h_i$.
\item Translate downward by $h_i e_3$ near $R_i'\times[0,h_i]$.
\end{enumerate}
A translation by $w$ is localized by
$\operatorname{curl}(\chi(w\times X)/2)$.
In the first and last quarters, choose cutoff collars below one third
of the smallest positive horizontal separation in the source and target
families, respectively. In the middle quarters, collars of thickness
one in height do not meet because the heights differ by four. For the
stretch, choose the horizontal plateau to include
$c_i+[-1,1]^2$. For translation, include the bounding rectangle of
$[c_i,c_i']$ enlarged by one in each horizontal direction.

More explicitly, with $\vartheta(s)=\sigma(2s-1/2)$ and
$\vartheta_j(s)=\vartheta(4s-j+1)$, multiply the sum of the stationary
fields for quarter $j$ by $\vartheta_j'(s)$. The tracked rectangle rises,
scales by $(\lambda_i^{\vartheta_2},\lambda_i^{-\vartheta_2})$,
translates its center along $[c_i,c_i']$, and descends. Each intermediate
half-width is at most $1/2$, so the entire rectangle remains on the
specified plateaus. The positive collars and finite affine norms also
give exactness on a positive initial neighborhood. If both endpoint
states are nonterminal, both first center coordinates are at most
$-3$, and the whole tracked path has negative first coordinate.

For the second realization choose $B=2m+1$, nonterminal offsets
$3,6,9,\ldots$, and terminal offsets $-3,-6,-9,\ldots$.
Thicken each rectangle by $[-1,1]$ in height, and write $c_i^0,c_i^1$
for the three-dimensional endpoint centers. Each pair within the
source family, and within the target family, has a horizontal coordinate
gap at least $B^{-2}$. Use height $6i$ and collar
$\epsilon=1/(4B^2)$. In three successive phases the prescribed affine
paths have center and linear part
\[
\begin{array}{c|c|c}
 &c_i(s)&L_i(s)\\\hline
 \text{lift}&c_i^0+6i\Theta e_3&I\\
 \text{transfer}&(1-\Theta)c_i^0+\Theta c_i^1+6i e_3&
     \operatorname{diag}(\lambda_i^\Theta,\lambda_i^{-\Theta},1)\\
 \text{lower}&c_i^1+6i(1-\Theta)e_3&
     \operatorname{diag}(\lambda_i,\lambda_i^{-1},1),
\end{array}
\]
where $\Theta$ is $\vartheta$ rescaled to the active third.
The proof of Lemma~\ref{lem:sb-boxes} applies to these exact data:
horizontal gaps separate the collars during lift and lowering, while
height gaps six separate vertical half-widths $1+\epsilon$ during
transfer. The localized field is explicitly
\begin{equation}\label{eq:p2:rapid-moving-box}
 \sum_i\operatorname{curl}_X\left[
 \chi_i\left\{\frac{\dot c_i\times(X-c_i)}2+
       \frac{(\dot L_iL_i^{-1}(X-c_i))\times(X-c_i)}3\right\}\right].
\end{equation}
The trace of $\dot L_iL_i^{-1}$ is zero. The curl identity in
Lemma~\ref{lem:sa-curl-motion} therefore gives the prescribed affine
velocity near each moving box, and ODE uniqueness gives its exact
full-box motion. A branch with two nonterminal endpoint states has
first center coordinate at least three and half-width at most $1/2$;
hence its entire path has first coordinate at least $5/2$.

It remains to initialize and to pass from the finite maps to the fluid
event. The initial combined tape is blank except at finitely many cells,
so its initial code $Z_0$ is rational. During $[0,1]$, a localized curl
translation with flat progress carries the fixed label to $Z_0$ along
the straight segment. A compact cutoff contains the whole segment with
a positive collar. Repeat the chosen instruction pulse on every later
unit interval. The time collars make all joins smooth, and only
finitely many spatial supports occur. Thus $U$ has the support and
derivative bounds in the statement. Set
\[
 f=\partial_tU+(U\cdot\nabla)U-\nu\Delta U.
\]
The residual realization and comparison argument of
Section~\ref{sec:model} gives the asserted solution and uniqueness;
fixed support and bounded $U$ give bounded kinetic energy. The formula
can also be returned as $f=f^{(0)}+\nu f^{(1)}$, with
$f^{(0)}=U_t+(U\cdot\nabla)U$ and $f^{(1)}=-\Delta U$ effective
independently of $\nu$. Evaluation for an arbitrary fixed real viscosity
is relative to that parameter.

At each integer $1+j$ through a halt, or at every such integer in a
nonhalting run, the particle is the exact code after $j$ auxiliary
instructions. Lemma~\ref{lem:sb-frontier} supplies finite positive
instruction counts between the original machine checkpoints. In the
first realization, a nonhalting loading segment has negative first
coordinate, and all subsequent nonterminal branches stay negative;
a terminal code has positive first coordinate. In the second,
nonhalting loading has nonnegative first coordinate and later
nonterminal paths stay above $5/2$ in that coordinate, while a terminal
code has first coordinate at most $-2$. An initially halted machine is
detected at the end of loading. Thus the exclusions apply at every real
time, and a halting run enters the detector no later than its terminal
checkpoint. No exclusion is needed during a final transition into a
terminal state.

All rectangles, cutoffs, scales and finite branch fields are computed
from the finite table. Positive rational scales have effective logarithms
and exponentials. Smooth switching and derivative bounds permit
evaluation at time joins without equality tests. The finite repeated
field acts on all branch domains at every period; neither its definition
nor its evaluation uses a predicted sequence of machine configurations.
\end{proof}
